\section{Historical Context and Literature Review}\label{sec:literature}

The historical measurement of capital accumulation in Latin America and Chile has evolved across three distinct traditions: regional quantitative economic history, official national accounts estimation, and political economy measurement debates regarding profit rates and capacity utilization.

\subsection{Regional Economic History and Latin American Capital Series}
The systematic reconstruction of historical capital formation in Latin America was inaugurated by \citet{Hofman2000a, Hofman2000b}. Using a standard perpetual inventory method (PIM) with fixed depreciation rates across broad asset categories, Hofman produced standardized gross and net capital stock estimates for six Latin American economies for 1950--1994, alongside investment expenditure series extending back to 1900. Hofman's contribution established the first benchmark for post-1950 capital levels in Chile, distinguishing machinery and equipment, non-residential construction, and residential housing. However, Hofman's primary objective was cross-country comparative growth accounting rather than enforcing strict annual stock--flow identities or deriving asset-specific scrappage behavior.

Building on the Latin American historical research program \citep{Carreras2006}, \citet{Tafunell2013} extended investment flows backward to 1856 across seventeen Latin American economies by constructing trade-based physical proxies (such as capital goods imports and apparent consumption of iron, steel, and cement). Subsequently, \citet{TafunellDucoing2016} produced long-run non-residential gross capital stock indices (1875--2008) for Argentina, Brazil, Chile, and Mexico. Their work emphasized physical productive capacity by deliberately excluding residential structures and applying variable service lives to equipment and structures. While \citet{TafunellDucoing2016} provided high-resolution relative growth profiles for non-residential gross stock ($K_g$), their series were published as normalized indices ($1929=100$) rather than real monetary levels, leaving the absolute level of installed capital unanchored.

\subsection{Chilean Historical National Accounts and Capital Stock Reconstructions}
At the national level, historical macroeconomic statistics for Chile center on the Pontificia Universidad Católica (PUC) Clio-Lab project, consolidated in \emph{La República en Cifras} \citep{Luders2016}. Within this framework, \citet{DiazWagner2016} estimated long-run net capital stocks ($K_n$) for Chile from 1833 to 2010 using standard PIM routines. The Clio-Lab series provided an invaluable quantitative backbone for Chilean economic history. However, raw Clio-Lab net series suffered from level scale distortions stemming from historical re-basing splices and focused exclusively on net book value ($K_n$), omitting gross capacity ($K_g$) and implicit retirement dynamics.

Earlier official PIM implementations by the Central Bank of Chile (\citealt{Aguilar2001}) provided net and gross capital estimates for 1985--2000 based on the 1993 System of National Accounts (SNA). Aguilar and Collinao demonstrated that PIM-based consumption of fixed capital frequently diverged from corporate tax depreciation, underscoring the sensitivity of net stock measures to accounting rules. More recently, \citet{BCCH2024} expanded official national accounts series to cover 1985--2024 in constant 2018 pesos (and historical 2003 CLP base surfaces), reporting real net stocks and depreciation by detailed asset class.

On the demand side, \citet{PerezEyzaguirre2019} undertook a major forensic revision of Chilean GDP and aggregate demand from 1860 to 1940. Pérez-Eyzaguirre corrected additivity and valuation flaws in earlier Clio-Lab gross fixed capital formation (GFCF) flows, producing the definitive constant 2003 CLP investment baseline. Pérez-Eyzaguirre's GFCF series provides the level-consistent flow foundation ($I_t$) for our harmonization.

\subsection{Capital Measurement Debates: Capacity, Wealth, and Stock--Flow Consistency}
In theoretical political economy and applied econometrics, the choice between gross and net capital measures is central to empirical modeling. Following \citet{Shaikh2016}, standard PIM implementations suffer from two key limitations: they treat depreciation rates as exogenous constants (e.g., fixed linear or geometric decay) and fail to distinguish between capital as physical productive capacity and capital as financial wealth.

Shaikh introduced the Generalized Perpetual Inventory Method (GPIM), demonstrating that gross capital ($K_g$) valued at replacement cost represents the appropriate base for measuring physical installed capacity and capital--output ratios ($\nu = K_g / Y$). Net capital ($K_n$), by contrast, incorporates cumulative accounting depreciation and is relevant for measuring net wealth, distributable surplus, and $r - g$ wealth dynamics \citep{Piketty2014}. As demonstrated by \citet{Basu2013} and \citet{Marquetti2010}, profit rate levels rotate significantly depending on whether historical-cost net stocks, replacement-cost net stocks, or gross physical stocks are employed.

Furthermore, \citet{Zarembka2015, Zarembka2019} emphasized that evaluating Marxian concepts such as the technical composition of capital ($K_g / L$) or the materialized composition of capital requires transparent separation between physical stock growth and price deflators. In applied macroeconomic research on Chile—such as recent investigations into long-run capacity utilization, profit rate dynamics, and accumulation regimes \citep{Duran2022, Polanco2019}—using unharmonized net series introduces severe specification bias.

Our paper bridges these literature strands by applying a Stock--Flow Consistent (SFC) harmonization framework to Chile. By unifying Hofman's level benchmarks, Pérez-Eyzaguirre's real investment flows, Tafunell \& Ducoing's sectoral gross indices, Clio-Lab's long-run growth profiles, and BCCh's modern national accounts, we construct a 125-year panel (1900--2024) that reconciles gross stocks, net stocks, investment, retirements, depreciation, and capital price deflators within a transparent, consistent framework.
